\section{The Atom}
    The atom is made up of a nucleus (a core of protons and neutrons).
    This nucleus makes up most of the mass of the atom, but is very small compared to the electron orbitals around it.
    Electrons are in `electron shells' around the nucleus, representing different energy values of the electrons.

    Protons are positively charged and electrons are negatively charged, therefore, for atoms to be neutral, there must be an equal number of protons and electrons.
    The mass of an atom is given (in relative atomic mass units) by the number of protons plus the number of neutrons, since the mass of electrons are negligable at this scale.
    An atom of element X can be written as its relative atomic mass (written A) and atomic number (written Z) (the number of protons).
    From its mass and atomic number, the number of neutrons can be found by subtracting the number of protons from the total mass.
    The number of electrons must be the same as the number of protons, since atoms are neutrally charged.
    The notation for the atom in this form is:
    \begin{equation}
        \ce{_Z^AX}
    \end{equation}
    For example, uranium-235 (a uranium atom (element number 92, symbol U) with a relative atomic mass of 235) would be written as:
    \begin{equation}
        \ce{^{235}_{92}U}
    \end{equation}
    \newpage

    \subsection{The Nucleus and Strong Nuclear Force}

\section{Radioactive Decay}
    \subsection{Alpha Decay}
        Alpha decay is where an unstable nucleus has to lose mass to become stable, so it releases an alpha particle ($\ce{^2_4\alpha}$), which is a particle made of 2 protons and 2 neutrons, equivalent to a helium-2 nucleus ($\ce{^4_2He^{2+}}$).
        This reduces the mass of the original atom by 4, and it loses 2 protons, changing its element.
        It will also leave the atom as an ion, which is called ionisation.
        Since the atom is losing 2 protons, the atom will have its atomic number decreased by 2.

        This is the equation for alpha decay of element X to element Y:
        \begin{equation}
            \ce{_Z^AX -> _{Z-4}^{A-2}Y + ^4_2\alpha}
        \end{equation}
        This can also be written with the alpha particle explicitly written as a helium nucleus:
        \begin{equation}
            \ce{_Z^AX -> _{Z-4}^{A-2}Y + ^4_2He}
        \end{equation}
        An example of this is the alpha decay of uranium-238 to thorium-234:
        \begin{equation}
            \ce{_92^235U -> _90^231Th + ^4_2\alpha}
        \end{equation}

    \subsection{Beta Minus Decay} \label{beta_minus}
        Beta minus decay is where a neutron turns into a proton, releasing an electron (also called a beta minus particle) and an electron antineutrino.
        Inside the decaying neutron, a down quark becomes an up quark, which releases the electron to conserve charge and the electron antineutrino to conserve lepton/electron number ($\leptonnumber[e]$).
        Since the atom gains a proton (and thus +1 charge), it becomes an ion, which is called ionisation.
        However, whereas alpha radiation decreases the atomic number by 2, beta minus radiation increases it by 1.

        The equation for beta minus decay of element X to element Y is:
        \begin{equation}
            \ce{_Z^AX -> _Z^{A+1}Y + _{-1}^0{\betaparticle[-]} + {\antineutrino[e]}}
        \end{equation}
        This can also be written with the beta particle explicity written as an electron:
        \begin{equation}
            \ce{_Z^AX -> _{Z+1}^AY + {\electron[-]} + {\antineutrino[e]}}
        \end{equation}
        An example of this is the beta minus decay of hydrogen-3 to helium-3:
        \begin{equation}
            \ce{_1^3H -> _2^3He + _{-1}^0{\betaparticle[-]} + {\antineutrino[e]}}
        \end{equation}

        The force responsible for beta minus decay is the weak nuclear force, mediated by the $\wboson[-]$ boson.

        Since a neutron is turning into a proton, the equation for beta minus decay can be written using the quark flavour of the neutron and proton, as $\ce{{\baryonquarks[udd]} -> {\baryonquarks[uud]} + {\betaparticle[-]} + {\antineutrino[e]}}$, which can be further simplified by removing the $\upquark$ and $\downquark$ that are left unchanged, to just leave $\ce{{\downquark} -> {\upquark} + {\betaparticle[-]} + {\antineutrino[e]}}$.

    \subsection{Beta Plus Decay}
        Beta plus decay is very similar to beta minus decay, however instead of a neutron becoming a proton, a proton becomes a neutron, releasing a positron (also called a beta plus particle), and an electron neutrino.
        Inside the decaying proton, an up quark becomes a down quark, which releases the positron to conserve the charge and the electron neutrino to conserve lepton/electron number ($\leptonnumber[e]$).
        Since the atom loses a proton (and thus -1 charge), it becomes an ion, which is called ionisation.
        However, whereas alpha radiation decreases the atomic number by 2, beta plus radiation decreases it by 1.
        The equation for beta plus decay of element X to element Y is:
        \begin{equation}
            \ce{_Z^AX -> _Z^{A-1}Y + _{1}^0{\betaparticle[+]} + {\neutrino[e]}}
        \end{equation}
        This can also be written with the beta particle explicity written as an electron:
        \begin{equation}
            \ce{_Z^AX -> _{Z-1}^AY + {\electron[+]} + {\neutrino[e]}}
        \end{equation}
        An example of this is the beta plus decay of hydrogen-3 to helium-3:
        \begin{equation}
            \ce{_{12}^{23}Mg -> _{11}^{23}He + _{1}^0{\betaparticle[+]} + {\neutrino[e]}}
        \end{equation}

        The force responsible for beta plus decay is the weak nuclear force, mediated by the $\wboson[+]$ boson.

        Since a proton is turning into a neutron, the equation for beta plus decay can be written using the quark flavour of the proton and neutron, as $\ce{{\baryonquarks[uud]} -> {\baryonquarks[udd]} + {\betaparticle[+]} + {\neutrino[e]}}$, which can be further simplified by removing the $\upquark$ and $\downquark$ that are left unchanged, to just leave $\ce{{\upquark} -> {\downquark} + {\betaparticle[+]} + {\neutrino[e]}}$.

        \newpage

\section{Feynman Diagrams}
    A Feynman diagram represents an interaction between particles.
    It has a time axis, showing the progression of the interaction from the incoming particles to the outgoing particles.
    
    This is a simple Feynman diagram, representing two electrons ($\electron[-]$) repelling each other.
    The interaction is mediated by the electromagnetic force, carried by a `virtual photon' ($\photon$):

    \begin{figure} \label{fig:electron_repulsion}
        \begin{center}
            \begin{tikzpicture}[x=15mm, y=15mm]
                \begin{feynman}
                    \vertex (i1) {e$^-$};
                    \vertex [above=1.5 of i1] (i2) {e$^-$};
                    
                    \vertex [right=2 of i1] (f1) {e$^-$};
                    \vertex [above=1.5 of f1] (f2) {e$^-$};

                    \vertex at ($(i1) + (.5,.75)$) (a);
                    \vertex at ($(f1) + (-.5,.75)$) (b);

                    \diagram*{
                        (i1) -- [fermion] (a) -- [photon, edge label=$\text{\gamma}$] (b),
                        (a) -- [fermion] (i2),
                        (f1) -- [fermion] (b),
                        (b) -- [fermion] (f2),
                    };
                    \draw[-stealth] (-.4,-.1) -- (-.4,1.35); % time axis
                    \node at (-.4,1.5) {$t$}; % time axis label
                \end{feynman}
            \end{tikzpicture}
        \end{center}
        \caption{A simple Feynman diagram.}
    \end{figure}

    Feynman diagrams can be used to show more complicated processes, such as beta minus decay ($\ce{{\downquark} -> {\upquark} + {\betaparticle[-]} + {\antineutrino[e]}}$):

    \begin{figure} \label{fig:beta_minus_decay}
        \begin{center}
            %https://tex.stackexchange.com/questions/617907/tikz-feynman-how-to-do-a-feynman-diagram-of-a-beta-decay
            \begin{tikzpicture}[x=15mm, y=15mm]
                \begin{feynman}
                    \vertex (i1) {u}; % up quark (left, n^0)
                    \vertex[right=.15 of i1] (i2) {d}; % down quark (middle, n^0)
                    \vertex[right=.15 of i2] (i3) {d}; % down quark (right, n^0)
                    \vertex[below=.3 of i2] (n) {n$^0$}; % neutron

                    \vertex[above=2 of i1] (f1) {u}; % up quark (left, p^+)
                    \vertex[right=.15 of f1] (f2) {d}; % down quark (middle, p^+)
                    \vertex[right=.15 of f2] (f3) {u}; % up quark (middle, p^+)
                    \vertex[above=.3 of f2] (p) {p$^+$}; % proton

                    \vertex[above=1 of i3] (a); % left u quark bend
                    \vertex[right=.15 of a] (b); % middle d quark bend
                    \vertex[right=.15 of b] (c); % right d->u quark bend

                    \vertex at ($(c) + (1,.25)$) (d); % left part of W boson line
                    \vertex at ($(d) + (.2, 1)$) (f4) {$\bar{\text{\nu}}_\text{e}$}; % electron antineutrino
                    \vertex at ($(d) + (.9,.5)$) (f5) {$\text{\beta}^-$}; % beta particle
                    
                    \diagram*{
                        (i1) -- [fermion] (a) -- [fermion] (f1), % left up quark
                        (i2) -- [fermion] (b) -- [fermion] (f2), % middle down quark
                        (i3) -- [fermion] (c) -- [fermion] (f3), % down->up quark
                        (c) -- [boson, edge label'=W$^-$] (d) -- [fermion] (f5), % W boson; beta particle
                        (d) -- [anti fermion] (f4), % electron antineutrino
                    };

                    \draw[-stealth] (-.4,-.4) -- (-.4,2.15); % time axis
                    \node at (-.4,2.3) {$t$}; % time axis label
                \end{feynman}
            \end{tikzpicture}
        \end{center}
        \caption{A Feynman diagram of $\betaparticle[-]$ decay; see section \ref{beta_minus}}
    \end{figure}

    In Feynman diagrams, fermions (leptons and baryons) are shown with solid straight lines with arrows representing their direction of travel through time.
    Matter travels forwards though time, and antimatter is modelled as travelling backwards, due to its charge being reversed from that of regular matter.
    Some diagrams show antimatter as travelling forwards in time along with regular matter.

    Most bosons (photons, $\wboson[\pm]$ bosons, $\zboson$ bosons) are shown with wavy lines (see the photon in figure \ref{fig:electron_repulsion} and the $\wboson[-]$ boson in \ref{fig:beta_minus_decay}).
    
    Mesons can be shown with a quark and antiquark line, or as a dashed line.

    \newpage